%%%%%%%%%%%%%%%%%%%%%%%%%%%%%%%%%%%%%%%%%%%%%%%%%%%%%%%%%%%%%%%%%%%%%%%%%%%%%%%%
% NeurIPS 2026 Paper: Ownership Verification for Diffusion Models
% via T-Error and Gaussian Quantile Regression

%%%%%%%%%%%%%%%%%%%%%%%%%%%%%%%%%%%%%%%%%%%%%%%%%%%%%%%%%%%%%%%%%%%%%%%%%%%%%%%%

\documentclass{article}
\usepackage{hyphenat}
% NeurIPS 2026 style — default (anonymous, double-blind) option
\usepackage{neurips_2026}

% ---------- Packages from NeurIPS template ----------
\usepackage[utf8]{inputenc}   % allow utf-8 input
\usepackage[T1]{fontenc}      % use 8-bit T1 fonts
\usepackage{hyperref}         % hyperlinks
\usepackage{url}              % simple URL typesetting
\usepackage{booktabs}         % professional-quality tables
\usepackage{amsfonts}         % blackboard math symbols
\usepackage{nicefrac}         % compact symbols for 1/2, etc.
\usepackage{microtype}        % microtypography
\usepackage{xcolor}           % colors

% ---------- Additional packages from ICML preamble ----------
\usepackage{amsmath,amssymb,amsthm}
\usepackage{graphicx}
\usepackage{subcaption}
\usepackage{multirow}
\usepackage{xspace}
\usepackage{pifont}
\usepackage{algorithm}
\usepackage{algorithmic}

% --- Figure (TikZ) ---
\usepackage{tikz}
\usetikzlibrary{positioning,arrows.meta,calc,fit,shapes.geometric}

% ---------- Custom commands ----------
\newcommand{\cmark}{\ding{51}}
\newcommand{\xmark}{\ding{55}}
\newcommand{\terr}{\ensuremath{e_t}\xspace}
\newcommand{\Wset}{\ensuremath{\mathcal{W}}\xspace}
\newcommand{\Dtrain}{\ensuremath{\mathcal{D}_{\text{train}}}\xspace}
\newcommand{\Dpublic}{\ensuremath{\mathcal{D}_{\text{public}}}\xspace}
\newcommand{\MIA}{\textsc{MIA}\xspace}
\newcommand{\OV}{\textsc{OV}\xspace}

% ---------- Theorem environments ----------
\newtheorem{definition}{Definition}
\newtheorem{claim}{Claim}

% ---------- Title ----------
\title{Membership is Ownership: A Robust Model Watermark Detection with Quantile Regression}

% ---------- Author block: anonymous for double-blind submission ----------
\author{%
  Anonymous Author(s)
}

\begin{document}

\maketitle

%%%%%%%%%%%%%%%%%%%%%%%%%%%%%%%%%%%%%%%%%%%%%%%%%%%%%%%%%%%%%%%%%%%%%%%%%%%%%%%%
% ABSTRACT
%%%%%%%%%%%%%%%%%%%%%%%%%%%%%%%%%%%%%%%%%%%%%%%%%%%%%%%%%%%%%%%%%%%%%%%%%%%%%%%%
\begin{abstract}
Large-scale diffusion models are valuable intellectual property whose weights are exposed to theft: an adversary who obtains a checkpoint can fine-tune and redeploy it without attribution. Existing approaches face structural limitations---training-time watermarking imposes a measurable utility cost and can be weakened by post-hoc fine-tuning, while membership inference attacks operate at the per-sample level and cannot anchor an ownership claim. We propose \textbf{Membership is Ownership} (MiO), an inference-time verification framework that reframes ownership verification as a population-level hypothesis-testing problem on a private member evidence set the owner curates and never releases. MiO comprises three components: (i) a multi-timestep $t$-error score aggregated via the lower quartile (\texttt{Q25}) to retain timesteps where the memorization gap is widest; (ii) a bagged ensemble of Gaussian quantile-regression networks that yields closed-form FPR thresholds at any target significance without per-quantile retraining; and (iii) a two-point verification protocol combining a consistency criterion (C1) and a multi-reference separation criterion (C2). Across four DDIM owners (CIFAR-10, CIFAR-100, STL-10, CelebA), every public reference yields paired separation $|d|_{\text{C2}} > 16$, an order of magnitude above the verification threshold, and the verdict persists under MMD fine-tuning, SGD fine-tuning, and weight-noise attacks. Across three Stable Diffusion v1.4 owner configurations, all (owner, adversary) cells are verified. To our knowledge, MiO is the first inference-time, model-ownership verification method for diffusion models; it matches the post-theft robustness of WDM, Zhao et al., and SleeperMark with zero weight modification, zero training overhead, and zero FID cost.
\end{abstract}

%%%%%%%%%%%%%%%%%%%%%%%%%%%%%%%%%%%%%%%%%%%%%%%%%%%%%%%%%%%%%%%%%%%%%%%%%%%%%%%%
% 1. INTRODUCTION
%%%%%%%%%%%%%%%%%%%%%%%%%%%%%%%%%%%%%%%%%%%%%%%%%%%%%%%%%%%%%%%%%%%%%%%%%%%%%%%%
\section{Introduction}
\label{sec:intro}

Large-scale generative diffusion models have become a central infrastructure for image synthesis~\citep{rombach2022high}. The capability of these models is bought at considerable expense, the high training cost makes the resulting weights valuable intellectual property (IP) for the institutions and individuals who train them. Yet these IP assets are exposed to a credible theft surface. An adversary who obtains a checkpoint can fine-tune and redeploy without attribution. Establishing whether a suspect checkpoint in the adversary's possession originates from a particular owner's base model—\emph{ownership verification}—is therefore a problem of practical economic and legal significance.

Existing approaches to diffusion-model IP protection address this problem along two complementary lines, each with structural limitations. \emph{Training-time watermarking} embeds an explicit signature into model parameters or output behavior~\citep{peng2023watermark,zhao2023recipe,wang2025sleepermark}: effective in controlled settings, but the embedded signature imposes a measurable utility cost and can be progressively weakened by post-hoc fine-tuning, which an adversary will always have means to apply. \emph{Membership inference attacks} (MIA) for diffusion models~\citep{kong2024efficient,matsumoto2023membership,hu2023loss} exploit the model's intrinsic memorization gap—reconstruction error on training samples is systematically lower than on unseen samples—and avoid weight modification entirely. However, MIA is fundamentally a per-sample classifier where a low reconstruction error on a single image cannot anchor an ownership claim, because individual samples may simply be easy to reconstruct.

The gap between these two lines suggests a third position. We propose \textbf{Membership is Ownership} (MiO), a verification framework built on a single principled observation: although the per-sample memorization signal is too weak to anchor ownership, the \emph{population-level} memorization pattern over a defender-curated evidence set is both robust and statistically auditable. Concretely, the owner reserves a private subset $\mathcal{W}$ of the training data; this set is never disclosed and serves as a non-invasive watermark that leaves model parameters and output quality entirely intact. To extract a robust signal from a suspect model, we introduce a multi-timestep \emph{$t$-error} score via forward noising and single-step reconstruction, and aggregate it via the lower quartile (\texttt{Q25}) to suppress noise while retaining the timesteps where the memorization gap is widest. To convert per-sample scores into calibrated decisions at arbitrary false-positive budgets, we fit a conditional Gaussian model in log-score space using a bagged ensemble of quantile-regression networks, yielding closed-form thresholds at any target FPR without per-quantile retraining.

This formulation moves ownership verification out of single-sample classification and into a population-level hypothesis-testing problem. Given an owner model $\mathcal{M}_A$ trained on data containing $\mathcal{W}$, a suspect model $\mathcal{M}_B$ submitted for verification, and a set of public reference checkpoints representing the null hypothesis, the verifier's task is to decide between \textsc{owned} and \textsc{independent} under a controlled false-positive rate. Our threat model assumes white-box access to both the owner model and the suspect model and permits the adversary to apply fine-tuning, distillation, or pruning to the stolen weights. Full retraining from scratch lies outside the threat model, as it severs provenance entirely.

Verification proceeds via two statistical criteria that test different nulls. \textbf{C1 (consistency)} requires that the owner and suspect models exhibit indistinguishable score \emph{distributions} on $\mathcal{W}$. \textbf{C2 (separation)} requires that the owner reconstructs $\mathcal{W}$ significantly better than every public reference. Ownership is declared verified only when both criteria hold jointly, ensuring that no single metric can yield a false positive.

We validate the framework across two evaluation scenarios chosen for complementary epistemic purposes. The first uses DDIM models trained from scratch on CIFAR-10, CIFAR-100, STL-10, and CelebA, providing the controlled environment necessary for measuring the protocol's quantitative properties with precision. The second evaluates Stable Diffusion v1.4 under LoRA , reflecting the production-scale threat model where actual diffusion-model IP theft occurs. Across all four DDIM datasets, every public reference yields paired separation $|d|_{\text{C2}} > 16$, an order of magnitude above the verification threshold; the separation persists under various post-theft adaptations, simulating realistic post-theft adaptation. Across all three SD owner configurations and both adversary types, the protocol returns a verified verdict in every case.

We summarize the contributions of this paper as follows.
\begin{itemize}
    \item \textbf{Conceptual.} We reframe ownership verification from single-sample membership inference into a population-level hypothesis-testing problem, using a private member evidence set as a non-invasive watermark.

    \item \textbf{Methodological.} We combine the multi-timestep \texttt{Q25} $t$-error score with a Gaussian quantile regression and a two-point verification protocol, yielding closed-form FPR control at any significance level without retraining.
    
    \item \textbf{Empirical.} We validate the framework on DDIM across four datasets 
    and on Stable Diffusion v1.4 across three owner configurations. All DDIM settings 
    achieve $|d|_{\text{C2}} > 16$ and remain verified under diverse post-theft 
    attacks; every SD configuration passes under adversarial fine-tuning.
    
    \item \textbf{Comparative.} Against WDM, Zhao et~al., CDI, and SleeperMark, 
    MiO matches post-theft robustness with zero weight modification, zero training 
    overhead, and zero FID cost.
\end{itemize}



%%%%%%%%%%%%%%%%%%%%%%%%%%%%%%%%%%%%%%%%%%%%%%%%%%%%%%%%%%%%%%%%%%%%%%%%%%%%%%%%
% 2. RELATED WORK
%%%%%%%%%%%%%%%%%%%%%%%%%%%%%%%%%%%%%%%%%%%%%%%%%%%%%%%%%%%%%%%%%%%%%%%%%%%%%%%%
\section{Related Work}
\label{sec:related}

\paragraph{Model Watermarking.}
Prior machine learning model watermarking methods typically embed an identifiable pattern into model behavior, most notably through trigger-based backdoors or carefully constructed inputs that induce a predetermined response~\citep{adi2018turning,zhang2018protecting,fernandez2023stable,wen2023tree}.
Such approaches provide a direct verification interface when the watermark remains intact, but they generally require explicit intervention during training and can be sensitive to subsequent model editing or ``cleaning.''
\textit{In contrast to embedding an external trigger, our setting treats a private member evidence set as the ``watermark'' and derives evidence from the model's intrinsic memorization behavior, avoiding modifications that introduce new behavioral artifacts.}

\paragraph{Membership Inference Attacks.}
Membership inference attacks (MIA) aim to decide whether an input was part of a model's training data~\citep{shokri2017membership,yeom2018privacy,salem2019ml}.
For diffusion models, recent work instantiates MIA using reconstruction-based signals~\citep{carlini2023extracting,duan2023diffusion,kong2024efficient}, loss/likelihood-based criteria~\citep{matsumoto2023membership,hu2023loss}, or scalable quantile-regression formulations~\citep{bertran2023scalable,tang2024membership}.
These methods are primarily evaluated as per-example decision rules (binary classification or scoring) and are typically reported via attack metrics.
Ownership verification, however, requires comparative, population-level evidence: the ownership claim must be supported by aggregate separation on a designated evidence set, calibrated against reference models that represent the null hypothesis, and summarized via statistically interpretable tests rather than single-sample predictions.
\textit{Our framework uses MIA-style signals as the underlying measurements, but creates an auditable verification procedure with public references and formal decision criteria.}



%%%%%%%%%%%%%%%%%%%%%%%%%%%%%%%%%%%%%%%%%%%%%%%%%%%%%%%%%%%%%%%%%%%%%%%%%%%%%%%%
% 3. PROBLEM SETUP
%%%%%%%%%%%%%%%%%%%%%%%%%%%%%%%%%%%%%%%%%%%%%%%%%%%%%%%%%%%%%%%%%%%%%%%%%%%%%%%%
\section{Problem Setup}
\label{sec:problem}

\subsection{Background and Problem Formulation}

\paragraph{Diffusion Models.}
A diffusion model corrupts a clean sample $x_0$ via a forward 
noising process,
\begin{equation}
    x_t = \sqrt{\bar\alpha_t}\, x_0 
        + \sqrt{1-\bar\alpha_t}\,\epsilon, 
    \quad \epsilon \sim \mathcal{N}(0, I),
\end{equation}
where $\bar\alpha_t = \prod_{s=1}^{t}\alpha_s$, and trains a 
network $\epsilon_\theta(x_t, t)$ to predict the added noise; 
sampling proceeds by iterative denoising.

Diffusion models systematically memorize portions of their training 
data~\citep{carlini2023extracting, somepalli2023diffusion}---a 
memorization gap our framework exploits in §\ref{sec:method}.

\paragraph{From Membership Inference to Ownership Verification.}
Membership inference attacks (MIAs) leverage the memorization gap
to determine whether a single query sample belongs to the training
set of a target model~\cite{duan2023diffusion, kong2024efficient}.
We state this formally:

\begin{definition}[Membership Inference Attack]
Given a target model $M$ and a query sample $x$, a
membership inference attack is a binary classifier
$\mathcal{A}: (M, x) \rightarrow \{0, 1\}$ that predicts
whether $x \in \mathcal{D}_{\mathrm{train}}$.
\end{definition}

However, a single low-error sample is insufficient evidence of 
ownership, since some samples are inherently easy to reconstruct. 
Ownership verification therefore requires a population-level test: 
significant memorization on a \emph{designated evidence set} 
relative to independent reference models.

\begin{definition}[Ownership Verification]
Let $\mathcal{M}_A$ be the owner's base model, $\mathcal{M}_B$ a 
suspect model, and $\mathcal{W} \subset \mathcal{D}_{\mathrm{train}}(\mathcal{M}_A)$ 
a private evidence set. An ownership verification protocol is a 
hypothesis test 
$\mathcal{V}: (\mathcal{M}_B, \mathcal{W}, \{M_{\text{ref}}^{(i)}\}) 
\rightarrow \{\texttt{owned}, \texttt{independent}\}$ 
that decides whether $\mathcal{M}_B$ derives from $\mathcal{M}_A$ 
under a controlled false-positive rate $\alpha$.
\end{definition}

\paragraph{Reference Model Selection.}
The reference set $\{M_{\text{ref}}^{(i)}\}_{i=1}^{|\mathcal{R}|}$ in the definition above is admissible only when each $M_{\text{ref}}^{(i)}$ satisfies three conditions: (i)~its training data is disjoint from $\mathcal{W}$;% supports: Phase 1 rows 5, 12; Phase 2 ID C1.1
(ii)~it has independent provenance from $\mathcal{M}_A$ and the owner;% supports: Phase 1 row 2; Phase 2 ID C2.1
and (iii)~it is type-compatible with the $t$-error scoring procedure of Section~\ref{sec:t_error}---i.e., a diffusion model on which a per-sample reconstruction score on $\mathcal{W}$ is well-defined.

\subsection{Threat Model}

We consider three principals: the \textbf{owner} $O$ who trained 
$\mathcal{M}_A$ and holds $\mathcal{W}$, the \textbf{adversary} 
who obtained $\mathcal{M}_A$ through theft and produced a 
derivative model $\mathcal{M}_B$ via fine-tuning, and the 
\textbf{verifier} $V$ who decides whether $\mathcal{M}_B$ derives 
from $\mathcal{M}_A$.

We assume the adversary has white-box access to $\mathcal{M}_A$
and may possess in-domain data disjoint from $\mathcal{W}$.The verifier 
has white-box access to both $\mathcal{M}_A$ and $\mathcal{M}_B$. The adversary is aware that a verification protocol exists but cannot
identify $\mathcal{W}$: even if $\mathcal{D}_{\mathrm{train}}$
were leaked, selecting the exact $|\mathcal{W}|$-subset from $N$
training samples is feasible with probability at most
%
\begin{equation}
    P(\text{exact match})
    = \frac{1}{\binom{N}{|\mathcal{W}|}}
    \;\leq\; \left(\frac{|\mathcal{W}|}{N}\right)^{|\mathcal{W}|},
\end{equation}
%
which is upper-bounded by $10^{-5000}$ under the CIFAR-10 setting
($N=50{,}000$, $|\mathcal{W}|=5{,}000$), making exhaustive guessing
computationally infeasible. 
The adversary is permitted to apply fine-tuning, distillation,
and pruning to $\mathcal{M}_A$; full retraining from scratch
falls outside our scope, as it severs provenance entirely.

Upon receiving a dispute, the verifier evaluates the protocol of 
Ownership Verification and outputs \textsc{Verified} or 
\textsc{Rejected} with supporting statistical evidence.


%%%%%%%%%%%%%%%%%%%%%%%%%%%%%%%%%%%%%%%%%%%%%%%%%%%%%%%%%%%%%%%%%%%%%%%%%%%%%%%%
% 4. METHOD
%%%%%%%%%%%%%%%%%%%%%%%%%%%%%%%%%%%%%%%%%%%%%%%%%%%%%%%%%%%%%%%%%%%%%%%%%%%%%%%%
\section{Method}
\label{sec:method}

The verification framework has four components:
(1) a \emph{t-error} score measuring reconstruction quality via forward-reverse diffusion,
(2) multi-timestep aggregation for stable per-sample scores, (3) Gaussian quantile regression for FPR-controlled thresholds,
and (4) a two-point statistical protocol comparing the suspect model against public baselines.

\subsection{$t$-error Scoring}
\label{sec:t_error}

Diffusion models reconstruct their training data with lower error than unseen data~\citep{carlini2023extracting}. We exploit this \emph{memorization gap} by measuring reconstruction quality at a given noise level $t$.

Specifically, given a clean sample $x_0$, we first corrupt it via the forward diffusion schedule:
\begin{equation}
x_t = \sqrt{\bar{\alpha}_t} \cdot x_0 + \sqrt{1 - \bar{\alpha}_t} \cdot \epsilon, \quad \epsilon \sim \mathcal{N}(0, I)
\end{equation}
where $x_t$ is the noised image at timestep $t \in \{1,\ldots,T\}$, $\epsilon$ is i.i.d.\ Gaussian noise, and $\bar{\alpha}_t = \prod_{s=1}^{t} \alpha_s$ is the cumulative noise coefficient under a cosine schedule~\citep{dhariwal2021diffusion}. The model then predicts the added noise $\hat{\epsilon}_\theta(x_t, t)$, from which we recover a single-step reconstruction:
\begin{equation}
\hat{x}_0 = \frac{x_t - \sqrt{1 - \bar{\alpha}_t} \cdot \hat{\epsilon}_\theta(x_t, t)}{\sqrt{\bar{\alpha}_t}}
\end{equation}
We define the \emph{$t$-error}, denoted $e_t(x_0)$, as the per-pixel mean squared error between the original and reconstructed images:
\begin{equation}
e_t(x_0) = \frac{\|x_0 - \hat{x}_0\|_2^2}{H \times W \times C}
\end{equation}
where $H$, $W$, $C$ denote the image height, width, and number of channels, respectively. The key property underlying our framework is the \emph{memorization effect}: models trained on a particular dataset exhibit systematically lower reconstruction error on their training examples, i.e., $\mathbb{E}_{x \in \mathcal{D}_{\mathrm{train}}}[e_t(x)] < \mathbb{E}_{x \notin \mathcal{D}_{\mathrm{train}}}[e_t(x)]$. This persistent gap, which survives across architectures and dataset scales, provides the foundational signal for our ownership verification protocol.

For latent diffusion models that operate in VAE latent space (e.g., Stable Diffusion), the definitions above transfer directly: the $t$-error computation is moved from pixel space to latent space, while the rest of the framework remains unchanged. The latent-vs-pixel design choice and its empirical justification are in Section~\ref{sec:ablation}.

\subsection{Multi-Timestep Aggregation}
\label{sec:aggregation}

A single timestep provides a noisy estimate of reconstruction quality, as the magnitude of $e_t(x)$ varies substantially with the noise level $t$. To obtain a stable per-sample statistic, we aggregate $t$-errors over $K$ uniformly sampled timesteps into a single score $s(x)$ via the 25th percentile:
\begin{equation}
  s(x) = Q_{0.25}\!\bigl(\{e_{t_i}(x)\}_{i=1}^{K}\bigr)
\end{equation}
We use $K{=}50$ by default. The choice of a lower percentile is deliberate: it captures the timesteps where the memorization gap is most pronounced while remaining robust to outlier timesteps that introduce high but uninformative error. Ablations in Section~\ref{sec:ablation} confirm that the 25th percentile yields the strongest separation among common aggregation strategies, and that performance is stable for $K \geq 25$. The full scoring procedure is summarized in Algorithm~\ref{alg:t_error}.

\begin{algorithm}[t]
\caption{$t$-error Computation}
\label{alg:t_error}
\begin{algorithmic}
\REQUIRE Model $\mathcal{M}$, sample $x_0$, timesteps $\mathcal{T} = \{t_1, \ldots, t_K\}$, noise schedule $\{\bar{\alpha}_t\}$
\ENSURE Aggregated score $s(x_0)$
\STATE $\texttt{errors} \leftarrow []$
\FOR{$t \in \mathcal{T}$}
    \STATE $\epsilon \sim \mathcal{N}(0, I)$
    \STATE $x_t \leftarrow \sqrt{\bar{\alpha}_t} \cdot x_0 + \sqrt{1 - \bar{\alpha}_t} \cdot \epsilon$
    \STATE $\hat{\epsilon} \leftarrow M_\theta(x_t, t)$
    \STATE $\hat{x}_0 \leftarrow (x_t - \sqrt{1 - \bar{\alpha}_t} \cdot \hat{\epsilon}) / \sqrt{\bar{\alpha}_t}$
    \STATE $e_t \leftarrow \|x_0 - \hat{x}_0\|_2^2 / (H \times W \times C)$
    \STATE $\texttt{errors}.\text{append}(e_t)$
\ENDFOR
\STATE \textbf{return} $s(x_0) \leftarrow Q_{0.25}(\texttt{errors})$
\end{algorithmic}
\end{algorithm}

\subsection{Two-Point Verification Criteria}
\label{sec:criteria}

Given the owner model $\mathcal{M}_A$, the suspect model 
$\mathcal{M}_B$, a set of public references $\{M_{\text{ref}}^{(i)}\}$, 
and the evidence set $\mathcal{W}$, we require both of the 
following criteria to be satisfied before declaring ownership 
verified.

\paragraph{Criterion 1: Consistency.}
If $\mathcal{M}_B$ derives from $\mathcal{M}_A$, the two models 
should exhibit similar reconstruction \emph{distributions} on 
$\mathcal{W}$. We test this with an unpaired Welch $t$-test on 
$\{s_{\mathcal{M}_A}(x)\}_{x\in\mathcal{W}}$ vs.\ 
$\{s_{\mathcal{M}_B}(x)\}_{x\in\mathcal{W}}$, requiring 
$p > 0.05$.

\paragraph{Criterion 2: Separation.}
The owner model must reconstruct watermark samples significantly 
better than any public reference. We apply C2 to the paired 
difference
\begin{equation}
\delta_{\text{C2}}(x) = s_{\mathcal{M}_A}(x) - s_{\text{ref}}(x)
\label{eq:delta_c2}
\end{equation}
via a one-sample $t$-test against zero, with effect size
\begin{equation}
|d|_{\text{C2}} = \frac{|\,\text{mean}(\delta_{\text{C2}}(\mathcal{W}))\,|}
{\text{std}(\delta_{\text{C2}}(\mathcal{W}))},
\label{eq:paired_c2}
\end{equation}
requiring both $p < 10^{-6}$ and $|d|_{\text{C2}} > 2.0$ 
(a "very large" effect~\citep{cohen1988statistical}).

Ownership is declared \textbf{verified} only when both criteria 
are jointly satisfied: C1 confirms provenance, while C2 establishes 
that the memorization signal is genuine rather than coincidental. 
The complete protocol is formalized in Algorithm~\ref{alg:verification}.

\begin{algorithm}[t]
\caption{Ownership Verification Protocol}
\label{alg:verification}
\begin{algorithmic}
\REQUIRE Models $M_A, M_B$, references $\{M_{\text{ref}}^{(i)}\}_{i=1}^{|\mathcal{R}|}$, evidence set $\mathcal{W}$
\ENSURE $\textsc{Verified}$ or $\textsc{Rejected}$
\STATE $S_A \leftarrow \{s_{M_A}(x) : x \in \mathcal{W}\}$; \quad $S_B \leftarrow \{s_{M_B}(x) : x \in \mathcal{W}\}$
\STATE \textbf{C1 (consistency, unpaired):} $p_1 \leftarrow \text{WelchTTest}(S_A, S_B)$
\IF{$p_1 < 0.05$} \STATE \textbf{return} $\textsc{Rejected}$ \ENDIF
\STATE \textbf{C2 (separation, paired):} \textbf{for each} reference $i \in \{1, \ldots, |\mathcal{R}|\}$:
\STATE \quad $S_{\text{ref}}^{(i)} \leftarrow \{s_{M_{\text{ref}}^{(i)}}(x) : x \in \mathcal{W}\}$
\STATE \quad $\delta_{\text{C2}}^{(i)} \leftarrow S_A - S_{\text{ref}}^{(i)}$ \COMMENT{element-wise on $\mathcal{W}$}
\STATE \quad $p_2^{(i)} \leftarrow \text{OneSampleTTest}\bigl(\delta_{\text{C2}}^{(i)},\, \mu_0 = 0\bigr)$
\STATE \quad $|d|_{\text{C2}}^{(i)} \leftarrow \bigl|\,\text{mean}(\delta_{\text{C2}}^{(i)})\,\bigr| \big/ \text{std}(\delta_{\text{C2}}^{(i)})$
\IF{$\exists i: p_2^{(i)} > 10^{-6}$ \textbf{or} $|d|_{\text{C2}}^{(i)} < 2.0$}
    \STATE \textbf{return} $\textsc{Rejected}$
\ENDIF
\STATE \textbf{return} $\textsc{Verified}$
\end{algorithmic}
\end{algorithm}

\subsection{Gaussian Quantile Regression}
\label{sec:gaussian_qr}

Rather than relying on fixed thresholds, we model the \emph{conditional score distribution} for non-members and derive decision boundaries in closed form. Concretely, we apply a log transformation to the aggregated score and parameterize the result as a Gaussian with input-dependent mean $\mu(x)$ and standard deviation $\sigma(x)$:
\begin{equation}
y = \log(1 + s(x)) \sim \mathcal{N}\!\bigl(\mu(x),\, \sigma^2(x)\bigr)
\end{equation}
where $\mu(x)$ and $\sigma(x)$ are jointly predicted by a neural network $f_\psi$. The network uses a ResNet-18 backbone, concatenated with a compact 
statistical summary $\phi(x) \in \mathbb{R}^3$ (mean, standard 
deviation, and $\ell_2$ norm of the $t$-error sequence), and an 
MLP head producing $(\mu, \log\sigma)$.

We train $f_\psi$ by minimizing the Gaussian negative log-likelihood 
over $N$ auxiliary non-member samples. Because the conditional distribution is Gaussian, the $(1{-}\alpha)$-quantile at any target false-positive rate $\alpha$ admits a closed-form expression:
\begin{equation}
\hat{q}_\tau(x) = \mu(x) + \sigma(x) \cdot \Phi^{-1}(\tau), \quad \tau = 1 - \alpha
\end{equation}
where $\Phi^{-1}$ is the standard normal inverse CDF. Unlike 
pinball-loss QR which requires a separate model per quantile 
level, a single $f_\psi$ supports arbitrary $\tau$.

The membership decision uses the \emph{margin}
\begin{equation}
m_\tau(x) = \hat{q}_\tau(x) - \log(1 + s(x)),
\end{equation}
where positive values indicate likely membership.

To reduce prediction variance, we wrap the procedure in a bagging 
ensemble of $B = 50$ models, each trained on a bootstrap sample 
of 80\% of the auxiliary data. The final ensemble margin 
$m_\tau^{\text{ens}}(x)$ averages individual quantile predictions; 
ensembling cost is discussed in Appendix~\ref{app:cost}.



%%%%%%%%%%%%%%%%%%%%%%%%%%%%%%%%%%%%%%%%%%%%%%%%%%%%%%%%%%%%%%%%%%%%%%%%%%%%%%%%
% 5. EXPERIMENTS
%%%%%%%%%%%%%%%%%%%%%%%%%%%%%%%%%%%%%%%%%%%%%%%%%%%%%%%%%%%%%%%%%%%%%%%%%%%%%%%%
\section{Experiments}
\label{sec:experiments}

\subsection{Experimental Setup}
\label{sec:setup}

\subsubsection{Evaluation Scenarios}
\label{sec:setup_scenarios}

We use two evaluation scenarios. DDIM trained from scratch on four datasets provides controlled measurement of the protocol's quantitative properties. Stable Diffusion v1.4 with PEFT instantiates the production-scale threat surface.

\subsubsection{Datasets and Models}
\label{sec:setup_datasets}

We evaluate our framework on four datasets spanning different resolutions and scales: CIFAR-10 and CIFAR-100 (both 32$\times$32, 50k training images with 5k reserved as the evidence set), STL-10 (96$\times$96, 5k training images with 1k as watermark), and CelebA (64$\times$64, 162k training images with 5k as watermark). All splits are deterministic and stratified with a fixed seed to ensure reproducibility.

For the Stable Diffusion experiments, the owner's private training set consists of 1{,}000 images from COCO 2014; per-image captions are retained for use in latent-space scoring (Section~\ref{sec:setup_training}). Adversaries fine-tune on a disjoint COCO subset (also 1{,}000 images) plus a synthetic-image set; details appear in Section~\ref{sec:setup_adversary}.

\subsubsection{Owner Training Protocols}
\label{sec:setup_training}

For each dataset, we train a DDIM model using a U-Net backbone with 128 base channels, a cosine noise schedule, and the AdamW optimizer at a learning rate of $2\times10^{-4}$ with EMA decay 0.9999. Training runs for 400k iterations in all cases. The Gaussian QR component uses an ensemble of $B{=}50$ models, each trained on a bootstrap sample of the auxiliary non-member data. Table~\ref{tab:hyperparams} summarizes the key hyperparameters; full architectural and optimization details are deferred to Appendix~\ref{app:details}.

\begin{table}[t]
\caption{Key hyperparameters for reproducibility. The DDIM and Gaussian QR rows are described in this subsection; the Stable Diffusion rows are split between owner training (described here) and adversary fine-tuning (described in Section~\ref{sec:setup_adversary} and Section~\ref{sec:theft}).}
\label{tab:hyperparams}
\centering
\begin{small}
\begin{tabular}{@{}lll@{}}
\toprule
Component & Parameter & Value \\
\midrule
\multirow{3}{*}{DDIM} & Iterations & 400k \\
& Batch size / LR & 128 / $2{\times}10^{-4}$ \\
& EMA decay & 0.9999 \\
\midrule
\multirow{3}{*}{Gaussian QR} & Ensemble $B$ & 50 \\
& Bootstrap ratio & 0.8 \\
& Epochs / patience & 50 / 10 \\
\bottomrule
\end{tabular}
\end{small}
\end{table}

For Stable Diffusion v1.4, owner models are produced by fine-tuning the base checkpoint on the COCO subset described in Section~\ref{sec:setup_datasets}. We consider three owner configurations: LoRA at rank $r{=}64$ (33.5M trainable parameters, 3.9\% of the UNet), LoRA at rank $r{=}256$ (134M, 15.6\%), and full UNet fine-tuning (860M, 100\%). Common settings are AdamW, cosine schedule with 500 warmup steps, batch size 4, fp16 mixed precision, seed 42; per-variant differences (steps, learning rate, augmentation) appear in Table~\ref{tab:hyperparams_sd} of Appendix~\ref{app:hyperparams}. The $t$-error is scored in VAE latent space rather than pixel space; the latent-vs-pixel ablation that motivates this choice is in Section~\ref{sec:ablation}.

\subsubsection{Adversary Protocols}
\label{sec:setup_adversary}

We consider four post-theft attacks on the DDIM owner. \textbf{MMD fine-tuning} runs 500 iterations of CLIP-feature distribution matching (cubic polynomial kernel over CLIP ViT-B/32 embeddings, $\mathcal{D}_{\text{train}}$ as the target distribution, lr $5{\times}10^{-6}$, EMA 0.9999) and represents an adversary attempting to alter the model's generation distribution while leaving the architecture intact. \textbf{SGD fine-tuning} (500 iterations, lr $5{\times}10^{-6}$, batch 128) on clean data represents a domain-preserving adversarial fine-tune. \textbf{Gaussian weight noise} at $\sigma \in \{10^{-3}, 10^{-2}\}$ applied in-place to all parameter tensors models a quantization or noise-perturbation adversary. The four attacks span the most plausible post-theft modifications under our threat model; quantitative results are in Section~\ref{sec:theft}.

For Stable Diffusion adversaries, a single 2{,}000-step LoRA fine-tune protocol (lr $5{\times}10^{-5}$, cross-attention modules only, $\epsilon$-prediction MSE loss) is instantiated across two axes. The data axis yields two threat scenarios: \textbf{domain-shift} on a disjoint 1{,}000-image COCO subset (the adversary adapts the stolen model to a different but distributionally similar dataset), and \textbf{task-shift} on synthetic data (the adversary repurposes the model toward a generation task different from the owner's). The rank axis varies the adversary's LoRA rank, $r \in \{64, 256\}$. The resulting $2{\times}2$ grid of (rank, scenario) adversary configurations is applied to each of the three owner configurations from Section~\ref{sec:setup_training}. Verification results across owners and scenarios are reported in Table~\ref{tab:sd_verification} of Section~\ref{sec:theft} (rank dimension collapsed); the SleeperMark comparison under the same adversary protocol is in Section~\ref{sec:baseline_comparison}.

\subsubsection{Baselines and Reference Models}
\label{sec:setup_baselines}

\paragraph{Public reference baselines.}
As public baselines for the verification protocol, we use pretrained checkpoints from HuggingFace: \texttt{google/ddpm-cifar10-32} serves as the reference for both CIFAR-10 and CIFAR-100, \texttt{google/ddpm-ema-bedroom-256} (resized) for STL-10, and \texttt{google/ddpm-celebahq-256} for CelebA. These baselines are trained on disjoint data from our evidence sets and represent the null hypothesis in the two-point verification protocol: a model that has never seen the owner's private evidence should exhibit substantially higher reconstruction error on that set.

\paragraph{Comparison methods.}
For the controlled baseline comparison in Section~\ref{sec:baseline_comparison}, we compare MiO against three diffusion-model IP protection methods: \textbf{WDM}~\citep{peng2023watermark}, a training-time watermarking method that learns a secondary watermark diffusion process; \textbf{Zhao et~al.}~\citep{zhao2023recipe}, which embeds StegaStamp fingerprints into training images before EDM training; and \textbf{CDI}~\citep{dubinski2024cdi}, an inference-time MIA-aggregation method. For Stable Diffusion, we additionally compare against \textbf{SleeperMark}~\citep{wang2025sleepermark}, a trigger-prompt-coupled, LoRA-resistant watermark. Implementation details and hyperparameters appear in Appendix~\ref{app:hyperparams}.

\subsubsection{Evaluation Metrics}
\label{sec:setup_metrics}
\paragraph{Verification verdicts.}
The two-point criteria of Section~\ref{sec:criteria} produce a binary $\{\textsc{Verified}, \textsc{Rejected}\}$ verdict per (owner, suspect) pair, together with the supporting C1 $p$-value and C2 effect size $|d|_{\text{C2}}$. Verification verdicts follow the C1/C2 criteria defined in Section~\ref{sec:criteria}.

\paragraph{Per-sample membership inference.}
To characterize the underlying signal, we report ROC-AUC and TPR at fixed FPR levels ($0.1\%$ and $0.01\%$) using the Gaussian QR ensemble margins of Section~\ref{sec:gaussian_qr}.

\paragraph{Effect sizes.}
Throughout, $|d|_{\text{C2}}$ denotes the paired Cohen's $d$ as defined in Equation~\ref{eq:paired_c2}. For per-sample membership analyses (e.g., owner-level AUC), we additionally report unpaired Cohen's $d$ between member and non-member score distributions, denoted $|d|_{\text{mem}}$.


\subsection{RQ1: Verification Effectiveness}
\label{sec:main_results}

\textbf{Can private member evidence sets distinguish owner models from public baselines?}

We test the two-point protocol on DDIM owners across four datasets, then validate the same protocol on Stable Diffusion. Three pieces of evidence support an affirmative answer.

\paragraph{C1 Consistency on DDIM.}
Owner $\mathcal{M}_A$ and suspect $\mathcal{M}_B$ produce 
statistically indistinguishable mean $t$-error on $\mathcal{W}$ 
across all four datasets: 
CIFAR-10: 28.73 vs.\ 28.74 ($p = 0.98$);
CIFAR-100: 33.35 vs.\ 33.34 ($p = 0.98$);
STL-10: 70.92 vs.\ 71.38 ($p = 0.71$);
CelebA: 63.79 vs.\ 63.80 ($p = 0.97$).
All $p$-values exceed the 0.05 acceptance threshold, confirming 
that $\mathcal{M}_B$ retains the memorization signature of its 
owner. C1 and C2 thus jointly hold across all four DDIM datasets.

\paragraph{C2 Separation on DDIM.}
Table~\ref{tab:main_results} reports Cohen's $d$ against three 
categories of public reference models (semantic-similar DDPM, 
cross-domain LSUN DDPM, and random initialization) under the 
ALL-must-pass criterion: every public reference must individually 
satisfy $|d| > 2.0$. The minimum $|d|$ across all references ranges 
from 16.3 (CIFAR-100, weakest dataset) to 22.2 (STL-10 with the 
LSUN-Church reference), exceeding the threshold by $8\times$ or 
more across all references. For STL-10, we use LSUN-Church as a 
conservative replacement for LSUN-Bedroom, which a prior reviewer 
flagged as a potential domain-match confound; the church reference 
yields $|d| = 22.2$, still $11\times$ above the $|d| > 2.0$ 
threshold. The two rightmost columns of Table~\ref{tab:main_results} 
report each dataset's minimum $|d|$ and its multiplicative margin; 
the smallest margin (CIFAR-100 at $8.1\times$) leaves substantial 
headroom, confirming that results are not artifacts of threshold 
tuning. These margins also confirm that the verification signal is 
an intrinsic property of the training data, not an artifact of a 
particular reference model choice.

\begin{table}[t]
\caption{Multi-reference ownership verification across 4 datasets.
Each owner model is tested against 3 public reference models grouped into
semantic-similar DDPM / cross-domain LSUN / random initialization.
Verification requires \emph{every} reference to satisfy $|d|>2.0$
(ALL-must-pass).
All $12$ (dataset, reference) pairs satisfy the threshold;
verification passes on all four datasets.
The rightmost two columns report the minimum $|d|$ achieved per dataset
and the multiplicative margin above the recommended threshold $|d|>2.0$.
All pairwise $p < 10^{-100}$.
$^\dagger$Conservative reference (LSUN-Church, replacing LSUN-Bedroom
to address domain-match concerns from prior review).}
\label{tab:main_results}
\centering
\begin{small}
\setlength{\tabcolsep}{4pt}
\begin{tabular}{@{}lccccc@{}}
\toprule
Dataset & Semantic-similar DDPM & Cross-domain LSUN DDPM & Random init & Min $|d|$ & Margin \\
        & $|d|$ & $|d|$ & $|d|$ & & \\
\midrule
CIFAR-10  & $20.7$ & $31.6$         & $45.1$ & $20.7$ & $10.4\times$ \\
CIFAR-100 & $16.3$ & $25.2$         & $44.1$ & $16.3$ & $8.1\times$ \\
STL-10    & $87.3$ & $22.2^\dagger$ & $41.5$ & $22.2$ & $11.1\times$ \\
CelebA    & $20.5$ & $23.5$         & $63.0$ & $20.5$ & $10.2\times$ \\
\bottomrule
\end{tabular}

\vspace{2pt}
\begin{minipage}{\linewidth}
\footnotesize
References per dataset: CIFAR-10/100 use
\texttt{google/ddpm-cifar10-32}, \texttt{google/ddpm-bedroom-256}, random-32.
STL-10 uses \texttt{google/ddpm-cifar10-32},
\texttt{google/ddpm-church-256}$^\dagger$, random-96.
CelebA uses \texttt{google/ddpm-celebahq-256},
\texttt{google/ddpm-bedroom-256}, random-64.
\end{minipage}
\end{small}
\end{table}





\paragraph{ Stable Diffusion Verification.}
We apply the verification protocol to the SD scenario described 
in Section~\ref{sec:setup_scenarios}, where the owner is SD~v1.4 
with private fine-tuning and the reference is the original SD~v1.4 
base. Table~\ref{tab:sd_verification} reports ownership verification 
across three owner configurations. The paired form of C2 in 
Section~\ref{sec:criteria} (Equation~\ref{eq:paired_c2}) is essential 
here: owner and reference share the SD base, so their raw $t$-errors 
are nearly identical, but $\delta_{\text{C2}}(x)$ cancels the shared 
per-image variance and exposes the systematic memorization gap. 
For reference, owner-level membership detection on $\mathcal{W}$ 
(1{,}000 members vs.\ 1{,}000 non-members, aggregated across both 
adversary fine-tuning directions) achieves AUC $= 0.995/0.999/0.999$ 
and $|d|_{\text{mem}} = 3.26/3.93/3.41$ for LoRA $r{=}64$ / LoRA 
$r{=}256$ / Full FT respectively. All three owner configurations 
satisfy both criteria across both adversary scenarios ($|d|_{\text{C2}} 
\in \{2.38, 2.43, 2.55\}$, $p_{\text{C1}} > 0.05$ uniformly), 
producing a \cmark\ verdict on all six (owner, adversary) cells of 
Table~\ref{tab:sd_verification}.

\begin{table}[t]
\caption{Ownership verification on SD~v1.4 (1{,}000 private COCO images).
C1: consistency ($p>0.05$).
C2: separation ($|d|>2.0$).
Adversaries: domain-shift fine-tuning on disjoint COCO and
task-shift on synthetic data.
C2 ($|d|$): effect size $|d|_{\text{C2}} = |\text{mean}(\delta_{\text{C2}}(\mathcal{W})) / \text{std}(\delta_{\text{C2}}(\mathcal{W}))|$ as defined in Equation~\ref{eq:paired_c2}.
$|d|_{\text{mem}}$: effect size between the owner's score
distributions on member and non-member samples in $\mathcal{W}$,
via standard absolute Cohen's $d$.}
\label{tab:sd_verification}
\centering
\small
\begin{tabular}{@{}llccccc@{}}
\toprule
Owner & Adversary & C1 & C2 ($|d|$) & AUC & $|d|_{\text{mem}}$ & Verdict \\
\midrule
LoRA $r\!=\!64$ & Domain-shift & \cmark~(p=0.73) & \cmark~(2.43) & 0.986 & 2.83 & \cmark \\
LoRA $r\!=\!64$ & Task-shift & \cmark~(p=0.34) & \cmark~(2.43) & 0.980 & 2.70 & \cmark \\
LoRA $r\!=\!256$ & Domain-shift & \cmark~(p=0.54) & \cmark~(2.55) & 0.998 & 3.51 & \cmark \\
LoRA $r\!=\!256$ & Task-shift & \cmark~(p=0.22) & \cmark~(2.55) & 0.997 & 3.46 & \cmark \\
Full FT & Domain-shift & \cmark~(p=0.71) & \cmark~(2.38) & 0.997 & 3.22 & \cmark \\
Full FT & Task-shift & \cmark~(p=0.18) & \cmark~(2.38) & 0.995 & 3.13 & \cmark \\
\bottomrule
\end{tabular}
\end{table}

\subsection{RQ2: Calibrated Per-Sample Detection at Stringent FPR Budgets}
\label{sec:fpr_control}

\textbf{Can per-sample membership scores be calibrated to the
stringent false-positive rates ownership disputes require?}

The Gaussian QR models the conditional log-score distribution for
non-members as a Gaussian with input-dependent mean and variance
$\mu(x), \sigma(x)$ (Section~\ref{sec:gaussian_qr}). The
$(1{-}\alpha)$-quantile
$\hat{q}_{1-\alpha}(x) = \mu(x) + \sigma(x)\,\Phi^{-1}(1{-}\alpha)$
gives a closed-form, per-sample threshold at any target FPR
$\alpha$---no retraining, no external calibration set, no
per-quantile tuning.

Table~\ref{tab:mia_results} reports the resulting per-sample AUC
together with TPR at $0.1\%$ and $0.01\%$ FPR on each owner model,
with $95\%$ bootstrap confidence intervals over 200 resamples.
Three facts from the table support the claim. \emph{First}, AUC
exceeds $0.96$ on all four datasets, confirming the underlying
$t$-error signal is strong. \emph{Second}, TPR remains above $0.55$
even at $0.01\%$ FPR (one false positive per ten thousand
non-members)---the regime ownership disputes operate in.
\emph{Third}, the same Gaussian QR ensemble produces useful
detection at \emph{both} FPR budgets without retraining, validating
the closed-form calibration empirically.

\begin{table}[t]
\caption{Per-sample MIA performance under analytical Gaussian QR
calibration, on each DDIM owner model. The same Gaussian QR
ensemble produces thresholds at both FPR budgets in closed form,
without retraining. TPR computed at fixed FPR with 95\% bootstrap
CI (200 iterations).}
\label{tab:mia_results}
\centering
\begin{small}
\begin{tabular}{lccc}
\toprule
Dataset & AUC & TPR@0.1\% FPR & TPR@0.01\% FPR \\
\midrule
CIFAR-10  & $0.987 \pm 0.002$ & $0.89 \pm 0.03$ & $0.71 \pm 0.05$ \\
CIFAR-100 & $0.982 \pm 0.003$ & $0.85 \pm 0.04$ & $0.65 \pm 0.06$ \\
STL-10    & $0.968 \pm 0.005$ & $0.78 \pm 0.05$ & $0.55 \pm 0.07$ \\
CelebA    & $0.975 \pm 0.004$ & $0.82 \pm 0.04$ & $0.60 \pm 0.06$ \\
\bottomrule
\end{tabular}
\end{small}
\end{table}
\subsection{RQ3: Post-Theft Robustness}
\label{sec:theft}

\textbf{Can MiO still detect ownership when the suspect model has 
undergone fine-tuning of unknown form?}

A practical adversary will not deploy the stolen $\mathcal{M}_A$ 
verbatim—they will fine-tune it to alter behavior, evade detection, 
or adapt to a new domain. The verification protocol must therefore 
remain effective across a diverse range of post-theft modifications, 
not just one. We test MiO under four distinct attack families on 
CIFAR-10:

\begin{itemize}
    \item \textbf{MMD fine-tuning}~\citep{gretton2012kernel}: 
    distribution-matching attack at the CLIP feature level
    \item \textbf{SGD fine-tuning}: standard gradient descent on 
    clean training data with the native $\varepsilon$-prediction loss
    \item \textbf{Gaussian weight noise} at $\sigma = 10^{-3}$ and 
    $\sigma = 10^{-2}$: in-place perturbation of all parameter tensors
\end{itemize}

These four attacks span the most plausible post-theft modifications: 
distribution-level adaptation, supervised fine-tuning, and weight 
perturbation. Detailed configurations appear in Appendix~\ref{app:hyperparams}.

\paragraph{MiO detects ownership across all four attack families.}
Table~\ref{tab:robustness} reports the verdict for each (method, 
attack) cell. MiO produces a \cmark\ verdict on every attack: the 
two-point criteria (C1, C2) remain jointly satisfied across MMD-FT, 
SGD-FT, and both noise levels. For MMD-FT specifically, Model~B 
retains nearly identical $t$-error to Model~A on $\mathcal{W}$ 
(Table~\ref{tab:main_results}), confirming that the memorization 
signal on individual evidence samples persists even when the model's 
generation distribution is being actively reshaped.

The mechanism behind this robustness is structural: because MiO 
operates on the model's intrinsic memorization of $\mathcal{W}$ 
rather than on an embedded behavioral signature, attacks that modify 
generation behavior (MMD-FT) or perturb weights moderately 
(noise, SGD-FT) do not erase the per-sample reconstruction gap on 
$\mathcal{W}$. Adversaries who retrain from scratch or specifically 
target memorization erasure fall outside our threat model.

\paragraph{Comparison with watermarking baselines.}
Table~\ref{tab:robustness} also includes WDM and Zhao under the 
same attack grid. Both baselines pass their respective thresholds 
across the cells where evaluation is feasible. MiO thus matches 
established watermarking methods on robustness while operating 
without their training-time cost or FID penalty (deferred analysis 
in Section~\ref{sec:baseline_comparison}).

\paragraph{Stable Diffusion.}
The SD verdict matrix in Table~\ref{tab:sd_verification} 
(§\ref{sec:main_results}) reports the same robustness pattern under 
the $2{\times}2$ (rank, scenario) adversary grid: all three owner 
configurations retain a \cmark\ verdict across both adversary 
scenarios, with $|d|_{\text{C2}} \in \{2.38, 2.43, 2.55\}$. Post-theft 
adaptation does not displace the memorization signal in either the 
DDIM-from-scratch or the SD-with-LoRA setting.
\begin{table}[t]
\caption{Robustness of watermark and ownership verification under
four post-theft attacks on CIFAR-10.  \cmark\ indicates the method
passes its native pass/fail threshold; ``---'' indicates not
applicable.  Pass thresholds: pixel accuracy $>0.9$ for WDM (on
the WDP-extracted watermark image); bit accuracy $>0.9$ for Zhao
(on $N\!=\!1000$ Heun-18 generated samples decoded with the
ground-truth 64-bit StegaStamp fingerprint); two-point verification
(C1: $p>0.05$, C2: $|d|>2.0$) for MiO.  WDM/Zhao thresholds are
MiO-defined; the original papers do not specify robustness
thresholds.  Attack configurations: MMD-FT and SGD-FT each run for
$500$ iterations at lr${}\!=\!5\!\times\!10^{-6}$ on clean CIFAR-10.
SGD-FT uses each method's native training objective
($\varepsilon$-prediction MSE for the DDPM-parameterized WDM and
MiO; $\sigma$-weighted $\hat{x}_0$-prediction MSE for Zhao's EDM
backbone) to avoid introducing a parameterization mismatch.
Gaussian weight noise is applied in-place to all parameter tensors
of each method's native network.
$^{\ddagger}$Zhao~$\times$~MMD-FT requires a Heun-based
differentiable sampler not implemented in our DDIM-based MMD
fine-tuning loop; deferred to future work.}
\label{tab:robustness}
\centering
\small
\begin{tabular}{@{}lcccc@{}}
\toprule
Method & MMD-FT & SGD-FT & $\sigma\!=\!10^{-3}$ & $\sigma\!=\!10^{-2}$ \\
\midrule
WDM            & \cmark & \cmark & \cmark & \cmark \\
Zhao$^\dagger$ & ---$^{\ddagger}$ & \cmark & \cmark & \cmark \\
\textbf{MiO}   & \cmark & \cmark & \cmark & \cmark \\
\bottomrule
\multicolumn{5}{@{}l}{\footnotesize $^\dagger$Uses EDM backbone (not DDPM/DDIM).} \\
\end{tabular}
\end{table}
\subsection{RQ4: Controlled Baseline Comparison}
\label{sec:baseline_comparison}

Section~\ref{sec:theft} established that MiO matches WDM and Zhao
on post-theft robustness. This section completes the comparison by
characterising MiO's position in the broader diffusion-IP landscape.
We benchmark MiO against the three CIFAR-10 baselines and the
SleeperMark SD baseline introduced in
Section~\ref{sec:setup_baselines}.\footnote{We exclude
DeepMarks~\citep{chen2019deepmarks}: its GMM-regularised weight
fingerprinting targets discriminative classifiers and has no public
implementation for generative architectures.}
WDM and Zhao are retrained on the same 50K CIFAR-10 images so that
any difference reflects the method itself rather than data or scale
effects; CDI is evaluated on the public
\texttt{google/ddpm-cifar10-32} checkpoint using the authors'
released codebase with minor compatibility patches
(Appendix~\ref{app:hyperparams}), under the default 7-attack
configuration with 5K members + 5K non-members.

\paragraph{Where MiO sits in the landscape.}
The four baselines differ from MiO in either \emph{what} they
verify or \emph{how} they verify it. WDM, Zhao, and SleeperMark
verify the same thing MiO does---that a suspect checkpoint
derives from the owner's model---but do so at training time, by
embedding a watermark. CDI works at inference time like MiO, but
verifies a different thing: whether the owner's \emph{data} was
used for training, not whether a model is the owner's.
MiO is, to our knowledge, the first method to verify model
ownership for diffusion models at inference time
(Table~\ref{tab:watermark_comparison}).

\paragraph{Capability cost vs. the watermarking baselines.}
WDM and Zhao both require a full (re)training pass with a modified
objective and incur measurable FID cost (+$7.21$ and +$3.07$
respectively, relative to their unwatermarked baselines under each
method's native sampler).  MiO modifies neither weights nor outputs.
Combined with the matched robustness from
Section~\ref{sec:theft}, the comparison shows that MiO attains the
same model-ownership guarantee at zero training and zero FID cost.

\paragraph{MIA signal strength vs. CDI.}
On the per-sample membership AUC axis, MiO scores $0.987$ and CDI $0.586$ on CIFAR-10.
This is not a head-to-head comparison because CDI is designed for
dataset-level aggregation in a weaker defender-information regime, and per-sample AUC understates its
intended capability. We report the gap to show that, when an
owner-side private evidence set is available, the per-sample MIA
signal alone is strong enough to drive verification.

\paragraph{Stable Diffusion (SleeperMark).}
Under the 2{,}000-step LoRA adversary protocol of
Section~\ref{sec:setup_adversary} (lr $5{\times}10^{-5}$,
cross-attention modules, 1{,}000 disjoint COCO images), MiO and
SleeperMark both pass their respective verification thresholds at
$r{=}64$ and $r{=}256$ (SleeperMark: bit-accuracy $> 0.9$; MiO: C1
+ C2 jointly satisfied).  MiO matches SleeperMark's verdict
without modifying the owner's weights or incurring training-time
watermarking cost.

\begin{table}[t]
\caption{Diffusion-model IP methods grouped by goal (model
vs.\ data ownership) and mechanism (training-time vs.\
inference-time).
FID cost is reported relative to each method's unwatermarked
baseline under its native sampler (WDM: 50-step DDIM; Zhao: 18-step
Heun on EDM).
Per-sample AUC is reported on CIFAR-10 for inference-time methods
only; per-sample classification is not the verification interface
of training-time watermarks.}
\label{tab:watermark_comparison}
\centering
\small
\setlength{\tabcolsep}{4pt}
\begin{tabular}{@{}llcccc@{}}
\toprule
Method & Goal & Mechanism & Training overhead & FID cost & Per-sample AUC \\
\midrule
WDM            & Model ownership & Training-time watermark & ${\sim}1\times$ & $+7.21$ & --- \\
Zhao$^\dagger$ & Model ownership & Training-time watermark & ${\sim}1\times$ & $+3.07$ & --- \\
\addlinespace
CDI            & \emph{Data} ownership   & Inference-time MIA & $0$           & $0$           & $0.586$ \\
\textbf{MiO}   & Model ownership & Inference-time MIA & $\mathbf{0}$  & $\mathbf{0}$  & $\mathbf{0.987}$ \\
\bottomrule
\multicolumn{6}{@{}l}{\footnotesize $^\dagger$Uses EDM backbone (not DDPM/DDIM).} \\
\end{tabular}
\end{table}
\subsection{Ablations}
\label{sec:ablation}

We validate four design choices on CIFAR-10 and on SD~v1.4 for 
the latent-vs-pixel choice; ensemble size and wall-clock 
breakdowns appear in Appendix~\ref{app:ablation}.

\begin{table}[t]
\caption{Aggregation ablation on CIFAR-10.  Lower percentiles
sharpen separation up to a point; Q25 offers the best trade-off
before outlier sensitivity sets in.
Ablation conducted on CIFAR-10; SD results use Q25 and are consistent with this ranking (see Section 5.5).}
\label{tab:aggregation}
\centering
\begin{small}
\begin{tabular}{lc}
\toprule
Aggregation & $|d|$ \\
\midrule
Mean & $18.2$ \\
Q10 & $21.5$ \\
Q25 (ours) & $\mathbf{23.93}$ \\
Median & $19.8$ \\
\bottomrule
\end{tabular}
\end{small}
\end{table}

\paragraph{Aggregation strategy.}
Table~\ref{tab:aggregation} compares four ways of summarising the 
$t$-error sequence. Q25 reaches $|d| = 23.93$, outperforming the 
mean (18.2), median (19.8), and the more aggressive Q10 (21.5); 
the latter begins to pick up outlier timesteps whose variance 
outweighs their selectivity.

\paragraph{Number of sampled timesteps.}
Varying $K$ from 10 to 100, $|d|$ plateaus once $K \geq 25$ and 
changes negligibly beyond $K = 50$. We fix $K = 50$ throughout; 
practitioners under tighter compute budgets can halve to $K = 25$ 
with only marginal loss.

\paragraph{Gaussian QR vs.\ pinball-loss QR.}
At $\tau \in \{0.999, 0.9999\}$, the two approaches yield 
comparable TPR. Gaussian QR's advantage is flexibility: one model 
supports any $\tau$ in closed form, whereas pinball-loss QR 
requires a separate model per quantile---a meaningful convenience 
for verifiers assessing claims at multiple significance levels.

\paragraph{SD: latent vs.\ pixel space.}
On SD~v1.4, latent-space scoring with per-image caption 
conditioning yields AUC $= 0.9956$, compared to $0.959$ for 
pixel-space scoring. Two factors explain the gap: the VAE decode 
step introduces a lossy reconstruction floor that dilutes the 
membership signal in pixel space, and empty prompts bypass the 
cross-attention pathway where LoRA adapters operate. All SD 
experiments use the latent-plus-caption configuration.
%%%%%%%%%%%%%%%%%%%%%%%%%%%%%%%%%%%%%%%%%%%%%%%%%%%%%%%%%%%%%%%%%%%%%%%%%%%%%%%%
% 6. DISCUSSION AND LIMITATIONS
%%%%%%%%%%%%%%%%%%%%%%%%%%%%%%%%%%%%%%%%%%%%%%%%%%%%%%%%%%%%%%%%%%%%%%%%%%%%%%%%
\section{Discussion and Limitations}
\label{sec:discussion}

\paragraph{The conceptual move.}
The contribution of this work is less a new diffusion-model
analysis tool than a relocation of where ownership evidence lives.
Embedded watermarks place the evidence \emph{inside the model}---
in parameters or outputs, where the owner can later recover an
artefact, and where an adversary can in principle attenuate it.
Our framework places the evidence \emph{outside the model}, in
the joint structure of the trained model and a private evidence
set the owner curates and never releases. The practical
consequence is that the protected asset is left intact: weights,
generation distribution, and FID are all unchanged. The
operational consequence is that the owner must maintain
$\mathcal{W}$ with auditable provenance, in the same way warranty
claims rely on retained receipts. Whether this trade-off is
preferable to embedded watermarking depends on how an institution
weighs output-quality cost against evidence-set maintenance cost;
we expect both approaches to coexist.

\paragraph{Scope of the verification claim.}
Three conditions delimit the regime in which our guarantees hold.
\emph{White-box access} to both the owner and suspect models is
required---the standard regime in IP litigation (courts compel
disclosure), regulatory audits (GDPR / AI~Act), partnership
disputes, and open-weight deployments.
\emph{Evidence-set retention with demonstrable provenance
predating the dispute}: training logs, dataset versioning, or
internal documentation. The combinatorial argument of
Section~\ref{sec:problem} bounds an adversary's ex-post guessing
probability of $\mathcal{W}$ at $\le 10^{-5000}$ under the
CIFAR-10 parameters, but the protocol does not by itself
adjudicate the temporal authenticity of $\mathcal{W}$; that
determination is delegated to the same evidentiary processes
governing other IP disputes.
\emph{Reference availability}: results are tightest with
domain-matched public references; for novel domains, general-purpose
references yield conservative rather than tight bounds.

\paragraph{Threat-model boundary.}
The verification signal rests on memorization that survives
weight-preserving adaptation (fine-tuning, pruning, quantisation).
Two attack classes lie outside this boundary. \emph{Knowledge
distillation} generates synthetic data from the stolen model and
trains a fresh student; the protocol correctly returns
\textsc{Fail} via C1 rejection
(Appendix~\ref{app:distillation}). This is a principled boundary
shared with embedded watermarking: WDM and Zhao both encode signals
into weights that a fresh-initialised student does not inherit.
\emph{Adaptive adversaries} who specifically optimise against the
memorization signal are not addressed here; the post-theft attacks
of Section~\ref{sec:theft} probe robustness against generic
adaptation, not against an adversary who knows the verification
mechanism and targets it. Characterising the resulting arms race
is open.

\paragraph{Open directions.}
Three directions remain open beyond the threat-model boundary.
First, a \emph{direct empirical FPR audit at the protocol level}---
running a battery of independently trained non-owner models
through the full C1+C2 pipeline and measuring per-trial rejection
rate---would upgrade the calibration analysis of
Section~\ref{sec:fpr_control} from analytical to empirical.
Second, \emph{API-only verification} is a substantively different
problem (the verifier loses direct access to the per-step
reconstruction signal) and is not addressed here.
Third, the \emph{relationship between memorization strength and
verifiability}---whether the protocol behaves differently on
heavily regularised, distilled-and-then-fine-tuned, or
explicitly de-memorised models---deserves systematic study, and
would clarify the boundary between "weight-preserving adaptation"
(handled) and "memorization-erasing adaptation" (not handled).
%%%%%%%%%%%%%%%%%%%%%%%%%%%%%%%%%%%%%%%%%%%%%%%%%%%%%%%%%%%%%%%%%%%%%%%%%%%%%%%%
% 7. CONCLUSION
%%%%%%%%%%%%%%%%%%%%%%%%%%%%%%%%%%%%%%%%%%%%%%%%%%%%%%%%%%%%%%%%%%%%%%%%%%%%%%%%
\section{Conclusion}
\label{sec:conclusion}

We reframed diffusion-model ownership verification as a
population-level hypothesis-testing problem on a private member
evidence set, and instantiated it through a multi-timestep
$t$-error score, Gaussian-QR-based analytical FPR calibration,
and a two-point (consistency, separation) verification protocol.
Across four DDIM datasets every public reference yields
$|d|_{\text{C2}} > 16$, an order of magnitude above the
verification threshold; on Stable Diffusion v1.4, all three owner
configurations remain verified under both LoRA-rank and adversary
scenarios.

To our knowledge this is the first inference-time, model-ownership
verification method for diffusion models---weights and outputs
remain unchanged, training is not modified, and FPR thresholds
are obtained in closed form. We hope the population-level framing
will prove useful beyond diffusion: any generative model with a
measurable per-sample memorization gap admits, in principle, the
same verification structure.

%%%%%%%%%%%%%%%%%%%%%%%%%%%%%%%%%%%%%%%%%%%%%%%%%%%%%%%%%%%%%%%%%%%%%%%%%%%%%%%%
% ACKNOWLEDGMENTS (hidden in anonymous mode by ack environment)
% Leave empty for anonymous submission. Add funding/thanks for camera-ready.
%%%%%%%%%%%%%%%%%%%%%%%%%%%%%%%%%%%%%%%%%%%%%%%%%%%%%%%%%%%%%%%%%%%%%%%%%%%%%%%%
\begin{ack}
\end{ack}

%%%%%%%%%%%%%%%%%%%%%%%%%%%%%%%%%%%%%%%%%%%%%%%%%%%%%%%%%%%%%%%%%%%%%%%%%%%%%%%%
% REFERENCES
%%%%%%%%%%%%%%%%%%%%%%%%%%%%%%%%%%%%%%%%%%%%%%%%%%%%%%%%%%%%%%%%%%%%%%%%%%%%%%%%
\bibliographystyle{plainnat}
\bibliography{references}


%%%%%%%%%%%%%%%%%%%%%%%%%%%%%%%%%%%%%%%%%%%%%%%%%%%%%%%%%%%%%%%%%%%%%%%%%%%%%%%%
% APPENDIX
%%%%%%%%%%%%%%%%%%%%%%%%%%%%%%%%%%%%%%%%%%%%%%%%%%%%%%%%%%%%%%%%%%%%%%%%%%%%%%%%
\appendix

\section{Additional Experimental Details}

\label{app:details}

\subsection{Hyperparameters}
\label{app:hyperparams}

\subsubsection{Owner Model Training}

\paragraph{DDIM owner training.} We train DDIM owner models for the four datasets in Section~\ref{sec:setup} using the configuration in Table~\ref{tab:hyperparams_ddim}.

\begin{table}[!htbp]
\caption{Training hyperparameters for DDIM models.}
\label{tab:hyperparams_ddim}
\centering
\begin{tabular}{ll}
\toprule
Parameter & Value \\
\midrule
Total iterations & 400,000 \\
Batch size & 128 \\
Learning rate & $2 \times 10^{-4}$ \\
Optimizer & AdamW \\
Weight decay & 0.0 \\
EMA decay & 0.9999 \\
Noise schedule & Cosine \\
Timesteps $T$ & 1000 \\
\midrule
\multicolumn{2}{c}{\textit{U-Net Architecture}} \\
Base channels & 128 \\
Channel multipliers & [1, 2, 2, 2] \\
Attention resolutions & [16] \\
Dropout & 0.1 \\
\bottomrule
\end{tabular}
\end{table}

\paragraph{Stable Diffusion fine-tuning.}
All three configurations share the same 1{,}000-image private subset of COCO 2014 as the evidence set $\mathcal{W}$; no per-configuration data resampling. We fine-tune SD~v1.4 from the official base checkpoint. Per-variant differences are summarized in Table~\ref{tab:hyperparams_sd}.

\begin{table}[!htbp]
\caption{Stable Diffusion owner fine-tuning configurations. Per-variant differences are listed above the rule; common settings shared by all three variants are listed below.}
\label{tab:hyperparams_sd}
\centering
\begin{small}
\begin{tabular}{@{}lccc@{}}
\toprule
Parameter & LoRA $r{=}64$ (A6) & LoRA $r{=}256$ (A8) & Full FT (A7) \\
\midrule
Trainable params       & 33.5M (3.9\%)     & 134M (15.6\%)     & 860M (100\%) \\
Steps                  & 20{,}000          & 20{,}000          & 10{,}000 \\
Learning rate          & $1{\times}10^{-4}$ & $1{\times}10^{-4}$ & $5{\times}10^{-6}$ \\
LoRA rank / $\alpha$   & 64 / 64           & 256 / 256         & --- \\
LoRA target modules    & cross-attn        & cross-attn        & --- \\
Random horizontal flip & yes               & no                & no \\
\midrule
\multicolumn{4}{c}{\textit{Common settings (all three variants)}} \\
\multicolumn{4}{l}{Base model: SD~v1.4 \quad Optimizer: AdamW \quad Batch size: 4 \quad Gradient accumulation: 1} \\
\multicolumn{4}{l}{Resolution: $512{\times}512$ resize with center crop \quad Captions: per-image COCO} \\
\multicolumn{4}{l}{LR schedule: cosine with 500 warmup steps \quad Mixed precision: fp16 \quad Seed: 42} \\
\bottomrule
\end{tabular}
\end{small}
\end{table}

\subsubsection{Gaussian QR Training}

\begin{table}[!htbp]
\caption{Gaussian QR training hyperparameters.}
\centering
\begin{tabular}{ll}
\toprule
Parameter & Value \\
\midrule
Backbone & ResNet-18 (pretrained) \\
Ensemble size $B$ & 50 \\
Bootstrap ratio & 0.8 \\
Epochs per model & 50 \\
Batch size & 256 \\
Learning rate & 0.001 \\
LR schedule & Cosine annealing \\
Early stopping patience & 10 epochs \\
\bottomrule
\end{tabular}
\end{table}

\subsubsection{Adversary Fine-tuning}

The adversary applies LoRA fine-tuning to either the owner's fine-tuned SD model (Table~\ref{tab:sd_verification}) or the SleeperMark watermarked checkpoint (Section~\ref{sec:baseline_comparison}). Identical settings in both cases: 2{,}000 steps, AdamW (weight decay $10^{-2}$), lr $5{\times}10^{-5}$, batch size 4, gradient clipping at norm 1.0, $\epsilon$-prediction MSE loss on 1{,}000 disjoint COCO images, seed 42. LoRA targets cross-attention modules \{\texttt{to\_q}, \texttt{to\_k}, \texttt{to\_v}, \texttt{to\_out.0}\} with rank/$\alpha \in \{64/64,\ 256/256\}$.

\subsubsection{Generative Baseline Training: WDM and Zhao}

Table~\ref{tab:baseline_training} summarizes the training configurations for the two generative-model baselines (WDM and Zhao). Method-specific watermark mechanisms are described after the table.

\begin{table}[!htbp]
\caption{Training configurations for generative baselines (CIFAR-10).}
\label{tab:baseline_training}
\centering
\begin{small}
\begin{tabular}{@{}lll@{}}
\toprule
Parameter & WDM & Zhao (EDM stage) \\
\midrule
Architecture            & U-Net 128ch                & SongUNet (ddpmpp) 128ch \\
Channel multipliers     & $[1,2,2,2]$                & $[2,2,2]$ \\
Residual blocks/stage   & 3                          & 4 \\
Attention resolutions   & $\{16, 8\}$, 4 heads       & $\{16\}$ \\
Activation / Norm       & SiLU / GroupNorm(32)       & (default EDM) \\
Dropout                 & 0.1                        & 0.13 \\
Noise schedule          & linear $\beta$, $T{=}1000$ & EDMPrecond \\
Prediction              & $\epsilon$                 & $\hat{x}_0$ ($\sigma$-weighted MSE) \\
Training duration       & 300{,}000 steps            & 15 Mimg (early-stopped)$^\dagger$ \\
Batch size              & 32                         & 512 (8 GPUs $\times$ 64) \\
Learning rate           & $1{\times}10^{-4}$         & $1{\times}10^{-3}$ \\
EMA decay               & 0.999                      & 0.5 \\
Optimizer               & AdamW                      & AdamW \\
Augmentation prob.      & ---                        & 0.12 \\
Seed                    & 42                         & 42 \\
FID sampling            & DDIM 50 steps              & Heun 18 steps \\
\bottomrule
\multicolumn{3}{l}{\footnotesize $^\dagger$Original Zhao paper trains to 200 Mimg; full duration would require approximately} \\
\multicolumn{3}{l}{\footnotesize 57 single-GPU-days. The 15 Mimg checkpoint achieves 100\% fingerprint bit accuracy,} \\
\multicolumn{3}{l}{\footnotesize sufficient for the baseline comparison.} \\
\end{tabular}
\end{small}
\end{table}

\paragraph{WDM watermark configuration.} Watermark blend factor $\gamma_1 = 0.8$, training loss weight $\gamma_2 = 0.1$ on watermarked samples ($\gamma_2 = 0.2$ at extraction), watermark seed 1998, fixed trigger image. Loss: $\mathcal{L} = \mathrm{MSE}(\epsilon) + \gamma_2 \cdot \mathrm{MSE}(\epsilon_{\text{wm}})$.

\paragraph{Zhao three-stage pipeline.} \textbf{Stage 1}: train StegaStamp encoder/decoder at $32{\times}32$ resolution, batch 64, 100 epochs, 64-bit fingerprint. \textbf{Stage 2}: embed a fixed 64-bit fingerprint (seed 42) into all 50K CIFAR-10 training images. \textbf{Stage 3}: train EDM with the configuration in Table~\ref{tab:baseline_training}.

\subsubsection{SleeperMark Baseline}

We use the pretrained SleeperMark U-Net checkpoint released by the authors (approximately 3.3 GB). No SleeperMark training is performed in our pipeline; the SleeperMark comparison in Section~\ref{sec:baseline_comparison} results measure robustness of the released checkpoint under the adversary protocol described above.

\subsubsection{CDI Reproduction}

We re-evaluated CDI~\citep{dubinski2024cdi} on \texttt{google/ddpm-cifar10-32} using the authors' codebase\footnote{\url{https://github.com/sprintml/copyrighted_data_identification}} on a single GPU (Python 3.8.18, PyTorch 2.4.1+cu121, diffusers 0.36.0; total wall-clock $\approx$49 minutes). \textbf{Configuration}: 5{,}000 members and 5{,}000 non-members sampled from CIFAR-10; batch size 100; seed 0; three-stage Hydra pipeline (features\_extraction $\to$ scores\_computation $\to$ evaluation). \textbf{Aggregation}: logistic regression (max\_iter 1000) over the union of four score-based attacks (denoising loss, SecMI-stat, PIA, PIAN) and three feature-based attacks (gradient masking, noise optimization, multiple-loss). Individual attack AUCs ranged from 0.491 (PIAN) to 0.557 (noise optimization); the aggregated CDI AUC reported in Table~\ref{tab:watermark_comparison} is 0.586.

\textbf{Patches.} Three trivial patches were required for the upstream \texttt{scores\_computation/cdi.py} aggregation stage to execute end-to-end: a missing \texttt{import numpy as np}; a \texttt{train\_dataset["data"]} $\to$ \texttt{train\_dataset.data} fix for the \texttt{MIDataset} class API; and a 3D$\to$2D reshape for multi-measurement features. None affects the aggregation algorithm itself.

\subsubsection{C2 Computation Details}

The Stable Diffusion verification numbers in Table~\ref{tab:sd_verification} are computed via Equation~\ref{eq:paired_c2}: for each adversary configuration, $\delta_{\text{C2}}(x) = s_{\text{tgt}}(x) - s_{\text{ref}}(x)$ is computed on $\mathcal{W}$ from the per-image score CSVs (one row per image, columns \texttt{score\_tgt} and \texttt{score\_ref}), and $|d|_{\text{C2}} = |\text{mean}(\delta_{\text{C2}}) / \text{std}(\delta_{\text{C2}})|$ is computed directly. A standalone runner that reproduces the Table~\ref{tab:sd_verification} C2 values from the released CSVs will be released with the camera-ready code.

\subsection{Post-theft Attack Configurations}

\paragraph{MMD fine-tuning.} 500 iterations on Model~A using a cubic polynomial kernel over CLIP ViT-B/32 (frozen) features. Optimizer: AdamW at lr $5{\times}10^{-6}$, batch size 128, EMA decay 0.9999. Generation pipeline: 10-step DDIM sampling from $t_{\text{start}}{=}900$. The MMD target distribution is $\mathcal{D}_{\text{train}}$ (which includes $\mathcal{W}$); see Section~\ref{sec:theft} for the rationale.

\paragraph{SGD fine-tuning.} 500 iterations, lr $5{\times}10^{-6}$, batch 128, applied on clean CIFAR-10 using each method's native training objective: $\epsilon$-prediction MSE for the DDPM-parameterized WDM and MiO; $\sigma$-weighted $\hat{x}_0$-prediction MSE for Zhao's EDM backbone. The choice of native objectives avoids introducing a parameterization mismatch.

\paragraph{Gaussian weight noise.} Additive isotropic Gaussian noise of standard deviation $\sigma \in \{10^{-3}, 10^{-2}\}$ applied in-place to all parameter tensors of each method's native network. No retraining; the perturbed checkpoint is evaluated directly.

\section{Gaussian QR Ensemble Size Ablation}
\label{app:ablation}

We vary the bagging ensemble size $B$ (introduced in Section~\ref{sec:gaussian_qr}) from 1 to 100 and measure the standard deviation of the predicted quantile $\hat{q}_\tau(x)$ across five independent runs on the CIFAR-10 evidence set at $\tau = 0.999$ (FPR $= 0.1\%$). Table~\ref{tab:ensemble_ablation} reports the results. Prediction variance drops sharply from $B = 1$ to $B = 20$, continues to decrease up to $B = 50$, and plateaus beyond that point. We adopt $B = 50$ as the default, which offers a favorable trade-off between stability and computational cost. Since ensemble members are trained independently, scaling from $B = 50$ to any larger value is embarrassingly parallel and requires no inter-model communication.

\begin{table}[!htbp]
\caption{Effect of ensemble size $B$ on quantile prediction stability (CIFAR-10, $\tau = 0.999$). Std.\ dev.\ computed over 5 independent runs.}
\label{tab:ensemble_ablation}
\centering
\begin{small}
\begin{tabular}{lcc}
\toprule
Ensemble $B$ & Std($\hat{q}_\tau$) & TPR @ FPR 0.1\% \\
\midrule
1 & 0.142 & 0.83 $\pm$ 0.05 \\
10 & 0.068 & 0.87 $\pm$ 0.03 \\
20 & 0.041 & 0.88 $\pm$ 0.03 \\
50 (ours) & 0.019 & 0.89 $\pm$ 0.03 \\
100 & 0.016 & 0.89 $\pm$ 0.02 \\
\bottomrule
\end{tabular}
\end{small}
\end{table}

\section{Computational Resources}
\label{app:cost}

Table~\ref{tab:cost} summarizes the computational requirements of each component. The dominant cost is ensemble training (approximately 100 A100-hours for $B = 50$), which parallelizes trivially. At inference, $t$-error computation for 5{,}000 watermark samples takes roughly 10 minutes on a single A100, and ensemble quantile prediction adds less than 1 minute. In total, a single ownership query (from raw images to verdict) completes in under 15 minutes, making the protocol practical for on-demand verification.

By comparison, explicit watermarking methods such as WDM~\citep{peng2023watermark} require re-training the diffusion model with the watermark objective, which typically takes as long as the original model training (on the order of days for the architectures we consider). MiO's Gaussian QR training is a one-time offline cost, and the owner model itself is never modified.

\begin{table}[!htbp]
\caption{Computational cost breakdown on CIFAR-10 (single NVIDIA A100).}
\label{tab:cost}
\centering
\begin{small}
\begin{tabular}{lcc}
\toprule
Component & Time & Parameters \\
\midrule
DDIM training (owner model) & $\sim$48 h & 35.7M \\
Gaussian QR ($\times 1$ model) & $\sim$2 h & 11.2M \\
Gaussian QR ensemble ($B{=}50$) & $\sim$100 h$^\dagger$ & 560M total \\
\midrule
\multicolumn{3}{l}{\textit{Inference (per verification query)}} \\
$t$-error scoring (5k samples) & $\sim$10 min & --- \\
Ensemble prediction & $<$1 min & --- \\
Two-point criteria & $<$1 s & --- \\
\midrule
\textbf{Total query time} & \textbf{$\sim$11 min} & --- \\
\bottomrule
\multicolumn{3}{l}{\footnotesize $^\dagger$ Embarrassingly parallel; wall-clock $\sim$2 h with 50 GPUs.} \\
\end{tabular}
\end{small}
\end{table}


\section{Gaussian QR Training Details}
\label{app:training}

\begin{algorithm}[!htbp]
\caption{Gaussian Quantile Regression Training}
\label{alg:gaussian_qr}
\begin{algorithmic}
\REQUIRE Auxiliary non-member data $\mathcal{D}_{\text{aux}}$, ensemble size $B$, bootstrap ratio $r$
\ENSURE Trained ensemble $\{f_{\theta}^{(b)}\}_{b=1}^{B}$
\FOR{$b = 1$ to $B$}
    \STATE $\mathcal{D}_b \leftarrow \text{BootstrapSample}(\mathcal{D}_{\text{aux}}, r)$ \COMMENT{80\% with replacement}
    \STATE $\mathcal{D}_{\text{train}}, \mathcal{D}_{\text{val}} \leftarrow \text{Split}(\mathcal{D}_b, 0.9)$
    \STATE Initialize $f_{\theta}^{(b)}$ with ResNet-18 backbone
    \REPEAT
        \FOR{batch $(x, s(x), \phi(x))$ in $\mathcal{D}_{\text{train}}$}
            \STATE $y \leftarrow \log(1 + s(x))$
            \STATE $(\mu, \log\sigma) \leftarrow f_{\theta}^{(b)}(x, \phi(x))$
            \STATE $\mathcal{L} \leftarrow (y - \mu)^2 / (2\sigma^2) + \log\sigma$
            \STATE Update $\theta$ via gradient descent
        \ENDFOR
    \UNTIL{early stopping (patience = 10)}
\ENDFOR
\STATE \textbf{return} $\{f_{\theta}^{(b)}\}_{b=1}^{B}$
\end{algorithmic}
\end{algorithm}


\section{Threat Model Boundary: Distillation}
\label{app:distillation}

An adversary who generates a large synthetic dataset from the stolen
model and trains a new model from scratch effectively performs
\emph{knowledge distillation}.  Because the student is initialised
independently and never sees the owner's training data directly,
per-example memorisation of $\mathcal{W}$ is not expected to transfer.

\paragraph{Experiment.}
We generate 50{,}000 synthetic CIFAR-10 images from Model~A via
10-step DDIM sampling and train a student model from scratch
(100{,}000 iterations, batch size 128, fresh random initialisation).
Table~\ref{tab:distillation} reports $t$-error statistics on $\mathcal{W}$.

\begin{table}[!htbp]
\caption{Distillation boundary experiment on CIFAR-10.
The student model, trained on synthetic data from the owner,
exhibits 3.72$\times$ higher $t$-error than the owner and fails
the consistency criterion.}
\label{tab:distillation}
\centering
\small
\begin{tabular}{@{}lccc@{}}
\toprule
Model & Mean $t$-error & Std & Ratio vs.\ Owner \\
\midrule
Model A (owner) & 28.71 & 10.85 & 1.0$\times$ \\
Student (distilled) & 106.90 & 36.14 & 3.72$\times$ \\
\addlinespace
ddpm-cifar10 (reference) & 703.84 & 38.50 & 24.5$\times$ \\
Random UNet (lower bound) & 3037.31 & 65.75 & 105.8$\times$ \\
\bottomrule
\end{tabular}
\end{table}

The student's mean $t$-error on $\mathcal{W}$ is $106.9$, compared to
$28.7$ for the owner—a $3.72\times$ gap. The verification protocol correctly returns \textsc{Fail}: the consistency criterion (C1) rejects because the student's score distribution differs significantly from the owner's ($p \approx 0$), indicating the student does not reproduce the owner's per-example memorisation pattern.

This outcome is expected and represents a principled boundary of the
method rather than a failure mode.  Distillation with fresh
initialisation approximates full retraining: generating a sufficiently
large synthetic dataset from a diffusion model and training a new model
from scratch approaches the computational cost of the original
training procedure, undermining the economic motivation for model theft.
Critically, the protocol does not produce a false positive---it
correctly reports that ownership cannot be established, which is the
far more important guarantee.

%%%%%%%%%%%%%%%%%%%%%%%%%%%%%%%%%%%%%%%%%%%%%%%%%%%%%%%%%%%%

\newpage
\input{checklist.tex}


\end{document}
